% -*- coding: utf-8 -*-
% !TEX program = xelatex
\documentclass[plain,noanswer,medmath]{../../template/jnuexam}
\usepackage{../../template/kysx}

\begin{document}


\examtitle{1991}{5}


\loadproblems{4}{1991P4}%载入试卷四


\makepart{填空题}{本题满分15分，每小题3分}


\useproblem{4}{1}{1}%【同试卷四第一(1)题】


\useproblem{4}{1}{2}%【同试卷四第一(2)题】


\useproblem{4}{1}{3}%【同试卷四第一(3)题】


\begin{problem}
$n$阶行列式$\begin{vmatrix}
       a &      b &      0 & \cdots &      0 & 0 \\
       0 &      a &      b & \cdots &      0 & 0 \\
       0 &      0 &      a & \cdots &      0 & 0 \\
  \vdots & \vdots & \vdots & \ddots & \vdots & \vdots \\
       0 &      0 &      0 & \cdots &      a & b \\
       0 &      0 &      0 & \cdots &      0 & a
\end{vmatrix} =$ \fillin{}．
\end{problem}


\begin{problem}
设$A$和$B$为随机事件，$P(A) = 0.7$, $P(A - B) = 0.3$，则$P(\overline{AB}) =$ \fillin{}．
\end{problem}


\makepart{选择题}{本题满分15分，每小题3分}


\useproblem{4}{2}{1}%【同试卷四第二(1)题】


\begin{problem}
设数列的通项为$x_n = \begin{cases}
  \frac{n^2 + \sqrt{n}}{n}, & \text{若$n$为奇数}; \\
               \frac{1}{n}, & \text{若$n$为偶数}.
\end{cases}$则当$n \to \infty$时，$x_n$是\pickout{}
\begin{abcd}
\item 无穷大量．   
\item 无穷小量．   
\item 有界变量．   
\item 无界变量．
\end{abcd}
\end{problem}


\begin{problem}
设$A$与$B$为$n$阶方阵，且$AB = O$，则必有\pickout{}
\begin{abcd}
\item $A = O$或$B = O$．					
\item $AB = BA$．
\item $|A| = 0$或$|B| = 0$．					
\item $|A| + |B| = 0$．
\end{abcd}
\end{problem}


\begin{problem}
设$A$是$m \times n$矩阵，$Ax = 0$是非齐次线性方程组$Ax = b$所对应的齐次线性方程组，
则下列结论正确的是\pickout{}
\begin{abcd}
\item 若$Ax = 0$仅有零解，则$Ax = b$有唯一解．
\item 若$Ax = 0$有非零解，则$Ax = b$有无穷多个解．
\item 若$Ax = b$有无穷多个解，则$Ax = 0$仅有零解．
\item 若$Ax = b$有无穷多个解，则$Ax = 0$有非零解．
\end{abcd}
\end{problem}


\useproblem{4}{2}{4}%【同试卷四第二(4)题】


\makepart{}{本题满分5分}

\begin{problem}
求极限$\lim_{x \to +\infty} (x + \sqrt{1 + x^2})^{\frac{1}{x}}$．
\end{problem}


\makepart{}{本题满分5分}

\begin{problem}
求定积分$I = \int_{-1}^1(2x + |x| + 1|)^2 \dx$．
\end{problem}


\makepart{}{本题满分5分}

\begin{problem}
求不定积分$I = \int \frac{x^2}{1 + x^2} \arctan x\dx$．
\end{problem}


\makepart{}{本题满分5分}

\begin{problem}
已知$xy = xf(z) + yg(z)$，$xf'(z) + yg'(z) \ne 0$，其中$z = z(x,y)$是$x$和$y$的函数．
求证：$[x - g(z)]\frac{\pdz}{\pdx} = [y - f(z)]\frac{\pdz}{\pdy}.$
\end{problem}


\makepart{}{本题满分6分}%【同试卷四第六题】

\useproblem{4}{6}{0}%


\makepart{}{本题满分8分}%【同试卷四第七题】

\useproblem{4}{7}{0}%


\makepart{}{本题满分6分}%【类似试卷四第八题】

\begin{problem}
证明不等式$\ln \left(1 + \frac{1}{x} \right) > \frac{1}{1 + x}$（$0 < x < +\infty$）．
\end{problem}


\makepart{}{本题满分5分}

\begin{problem}
设$n$阶矩阵$A$和$B$满足条件$A + B = AB$．
\begin{enumerate}
\item 证明$A - E$为可逆矩阵，其中$E$是$n$阶单位矩阵；
\item 已知$B = \begin{pmatrix}
  1 & -3 & 0 \\
  2 &  1 & 0 \\
  0 &  0 & 2
\end{pmatrix}$，求矩阵$A$．
\end{enumerate}
\end{problem}


\makepart{}{本题满分7分}%【同试卷四第九题】

\useproblem{4}{9}{0}%


\makepart{}{本题满分5分}

\begin{problem}
已知向量$\alpha =(1,k,1)^T$是矩阵$A = \begin{pmatrix}
  2 & 1 & 1 \\
  1 & 2 & 1 \\
  1 & 1 & 2
\end{pmatrix}$的逆矩阵$A^{-1}$的特征向量，求常数$k$的值．
\end{problem}


\makepart{}{本题满分7分}%【类似试卷四第十二题】

\begin{problem}
一汽车沿一街道行驶，需要通过三个均设有红绿信号灯的路口，
每个信号灯为红或绿与其他信号灯为红或绿相互独立，且红绿两种信号显示的时间相等，
以$X$表示该汽车首次遇到红灯前已通过的路口的个数．\par
\begin{enumerate*}
\item 求$X$的概率分布．
\item 求$E\left(\frac{1}{1 + X} \right)$．
\end{enumerate*}
\end{problem}


\makepart{}{本题满分6分}

\begin{problem}
在电源电压不超过$200$伏、在$200 - 240$伏和超过$240$伏三种情形下，
某种电子元件损坏的概率分别为$0.1$,$0.001$和$0.2$，
假设电源电压$X$服从正态分布$N(220,25^2)$，试求：
\begin{enumerate}
\item 该电子元件损坏的概率$\alpha$；
\item 该电子元件损坏时，电源电压在$200 - 240$伏的概率$\beta$．
\end{enumerate}
\centerline{[附表]（表中$\Phi(x)$是标准正态分布函数）}\noindent
$$\begin{array}{|c|*{8}{c}|}
\hline
        x &  0.10 &  0.20 &  0.40 &  0.60 &  0.80 &  1.00 &  1.20 & 1.40 \\
\hline
  \Phi(x) & 0.530 & 0.579 & 0.655 & 0.726 & 0.788 & 0.841 & 0.885 & 0.919 \\
\hline
\end{array}$$
\end{problem}


\end{document}
